\documentclass[11pt]{article}
\usepackage[margin=1in]{geometry}
\usepackage{amsmath,amssymb,amsthm}
\usepackage{graphicx}
\usepackage[hidelinks]{hyperref}

\newtheorem{theorem}{Theorem}
\newtheorem{corollary}{Corollary}
\newtheorem{lemma}{Lemma}
\newtheorem{remark}{Remark}

\newcommand{\R}{\mathbb R}
\newcommand{\M}{\mathbb M}
\newcommand{\MU}{\operatorname{MU}}
\newcommand{\Ad}{\operatorname{Ad}}
\newcommand{\ceil}[1]{\left\lceil #1\right\rceil}

\title{Finite-Component Structure of Mixed-Unitary Times}
\author{FA Banach run \texttt{fa\_banach\_001}, agent \texttt{agent\_lane\_10}}
\date{23 July 2026}

\begin{document}
\maketitle

\section*{Status}

This packet records a candidate partial result, likely valid, toward the open
remark following Theorem 5.4 of Bhat and Devendra \cite{BD}. The source asks
whether the first mixed-unitary re-entry time must already be the beginning of
the eventual mixed-unitary ray. We do not decide that existence question.
Instead, we prove that the full set of mixed-unitary times always has only
finitely many connected components.

\section*{Source Passage}

\begin{center}
\includegraphics[width=0.98\linewidth]{figures/open_problem_crop.png}
\end{center}

\section*{Question and Answered Part}

Let \((\Phi_t)_{t\ge0}\), with \(\Phi_t=e^{t\mathcal L}\), be a quantum
dynamical semigroup of unital channels on \(\M_d\). The source proves that
either every \(\Phi_t\) is mixed unitary, or there are
\(0<t_0\le t_1<\infty\) such that
\[
 \Phi_t\notin\MU(d)\quad(0<t<t_0),\qquad
 \Phi_{t_0}\in\MU(d),\qquad
 \Phi_t\in\MU(d)\quad(t\ge t_1).
 \tag{1}\label{eq:source}
\]
It asks whether \(t_0=t_1\) must hold. If not, the authors suggest that the
mixed-unitary times in \([t_0,t_1]\) might form a discrete set.

We prove that arbitrary infinite oscillation cannot occur: the complete time
set is a finite union of intervals and points. In particular, a discrete
intermediate set must be finite.

\section*{Proof Intuition}

The set of mixed-unitary channels in fixed dimension is semialgebraic. Indeed,
the unitary group is cut out by polynomial equations in real and imaginary
parts, unitary conjugation depends polynomially on those coordinates, and
Carath\'eodory's theorem gives a uniform finite bound on the number of terms
in a convex combination. Tarski--Seidenberg then removes the auxiliary
unitaries and convex coefficients.

A quantifier-free semialgebraic description uses the signs of only finitely
many polynomials. On the analytic curve \(t\mapsto e^{t\mathcal L}\), those
polynomials become finitely many real-analytic functions. A nonzero
real-analytic function has only finitely many zeros on a compact interval.
Consequently the membership pattern can change only finitely many times.

\section*{Semialgebraicity}

We regard the real vector space of Hermiticity-preserving linear maps
\(\M_d\to\M_d\) as \(\R^N\), after fixing real coordinates.

\begin{lemma}\label{lem:semi}
For every \(d\), the set \(\MU(d)\) of mixed-unitary channels on \(\M_d\) is
a compact semialgebraic subset of \(\R^N\).
\end{lemma}

\begin{proof}
Write a complex \(d\times d\) matrix in real and imaginary coordinates.
The unitary group \(U(d)\) is semialgebraic because \(U^*U=I\) is a finite
system of real polynomial equations. In the fixed coordinates on the space of
linear maps, the entries of
\[
 \Ad_U:X\longmapsto UXU^*
\]
are real polynomials in the entries of \(U\).

By Carath\'eodory's theorem, every point in the convex hull of a subset of
\(\R^N\) is a convex combination of at most \(N+1\) points. Hence
\(\Phi\in\MU(d)\) if and only if there exist
\[
 U_0,\ldots,U_N\in U(d),\qquad
 \lambda_0,\ldots,\lambda_N\ge0,\qquad
 \sum_{j=0}^N\lambda_j=1,
\]
such that
\[
 \Phi=\sum_{j=0}^N\lambda_j\Ad_{U_j}.
 \tag{2}\label{eq:convex}
\]
All conditions in \eqref{eq:convex} are polynomial equalities or
inequalities in finitely many real variables. Thus the set of tuples
\((\Phi,\lambda_0,\ldots,\lambda_N,U_0,\ldots,U_N)\) satisfying them is
semialgebraic. Its projection to the \(\Phi\)-coordinate is semialgebraic by
the Tarski--Seidenberg theorem \cite[Chapter 2]{BPR}.

Finally, \(U(d)\) is compact, and the convex hull of a compact subset of a
finite-dimensional space is compact. Therefore \(\MU(d)\) is compact.
\end{proof}

\section*{Finite-Component Theorem}

\begin{theorem}\label{thm:main}
Let \((\Phi_t)_{t\ge0}\), \(\Phi_t=e^{t\mathcal L}\), be a quantum
dynamical semigroup of unital channels on \(\M_d\), and define
\[
 S_{\mathcal L}
 =
 \{t\ge0:\Phi_t\in\MU(d)\}.
 \tag{3}\label{eq:times}
\]
Then:
\begin{enumerate}
\item \(S_{\mathcal L}\) is a closed additive subsemigroup of
      \([0,\infty)\) containing \(0\);
\item \(S_{\mathcal L}\) has finitely many connected components;
\item more explicitly, \(S_{\mathcal L}\) is a finite union of isolated
      points and closed intervals, one interval being a ray
      \([\tau,\infty)\).
\end{enumerate}
\end{theorem}

\begin{proof}
The identity channel \(\Phi_0\) is unitary, so \(0\in S_{\mathcal L}\).
The set is closed because \(t\mapsto\Phi_t\) is continuous and \(\MU(d)\) is
closed. If \(s,t\in S_{\mathcal L}\), then
\(\Phi_{s+t}=\Phi_s\circ\Phi_t\) is mixed unitary: composing two convex
combinations of unitary conjugations gives another such combination.
Therefore \(s+t\in S_{\mathcal L}\).

By eventual mixed unitarity, Theorem 4.12 of \cite{BD}, there exists
\(T<\infty\) such that
\[
 [T,\infty)\subseteq S_{\mathcal L}.
 \tag{4}\label{eq:eventual}
\]
It remains to prove that \(S_{\mathcal L}\cap[0,T]\) has finitely many
components.

By Lemma \ref{lem:semi} and quantifier elimination for semialgebraic sets,
membership in \(\MU(d)\) has a quantifier-free description involving the
signs of finitely many real polynomials
\[
 P_1,\ldots,P_m
\]
in the coordinates of a linear map. Put
\[
 f_j(t)=P_j(\Phi_t)=P_j(e^{t\mathcal L}).
\]
Every coordinate of the finite-dimensional matrix exponential
\(e^{t\mathcal L}\) is real analytic in \(t\). Hence every \(f_j\) is real
analytic on \(\R\).

If \(f_j\) is not identically zero, it has only finitely many zeros in
\([0,T]\). Indeed, infinitely many zeros in that compact interval would have
an accumulation point, forcing a real-analytic function to vanish
identically. Functions that are identically zero impose a constant sign
condition. Let \(Z\) be the union of the zero sets in \([0,T]\) of all
nonzero \(f_j\). Then \(Z\) is finite.

On every component of \([0,T]\setminus Z\), the sign of every \(f_j\) is
constant. The quantifier-free membership formula therefore has a constant
truth value on that component. It follows that
\(S_{\mathcal L}\cap[0,T]\) is a finite union of intervals and points.
Because \(S_{\mathcal L}\) is closed, the intervals may be taken closed
after merging adjacent pieces. Combining them with the ray
\eqref{eq:eventual} proves the result.
\end{proof}

\begin{corollary}[The discrete scenario is finite]\label{cor:discrete}
Assume the nontrivial alternative \eqref{eq:source} from the source paper.
Then the mixed-unitary times in \([t_0,t_1]\) form a finite union of closed
intervals and isolated points. In particular, if that set is discrete, it is
finite.
\end{corollary}

\begin{proof}
Apply Theorem \ref{thm:main}. A closed discrete subset of a compact interval
is finite.
\end{proof}

\begin{corollary}[An interval forces a ray]\label{cor:ray}
If \(0<a<b\) and \([a,b]\subseteq S_{\mathcal L}\), then, for
\[
 n_0=\ceil{\frac{a}{b-a}},
\]
one has
\[
 [n_0a,\infty)\subseteq S_{\mathcal L}.
 \tag{5}\label{eq:ray}
\]
In particular, an interval \([0,b]\subseteq S_{\mathcal L}\) forces
\(S_{\mathcal L}=[0,\infty)\).
\end{corollary}

\begin{proof}
Additivity gives
\[
 m[a,b]=[ma,mb]\subseteq S_{\mathcal L}
 \qquad(m\in\mathbb N).
\]
For every \(m\ge n_0\),
\[
 mb\ge(m+1)a
 \quad\Longleftrightarrow\quad
 m(b-a)\ge a.
\]
Thus the intervals \([ma,mb]\), \(m\ge n_0\), overlap consecutively, and
their union is \([n_0a,\infty)\). If \([0,b]\subseteq S_{\mathcal L}\),
every \(t\ge0\) is a sum of finitely many elements of that interval.
\end{proof}

\section*{Limitations}

Theorem \ref{thm:main} rules out a Cantor-like time set and rules out
infinitely many isolated re-entries before eventual mixed unitarity. It does
not prove \(t_0=t_1\). For example, at the purely additive level, a set such
as
\[
 \{0,a\}\cup[2a,\infty)
\]
is a closed additive subsemigroup with a genuine gap. The remaining problem
is whether the analytic curve generated by a Lindbladian can realize such a
pattern inside the particular convex body \(\MU(d)\).

\section*{Novelty Check and Human Review}

The run indexes and bounded external searches on 23 July 2026 were checked
using the arXiv id, exact title, mixed-unitary time sets, semialgebraicity,
finite unions of intervals, and semigroup formulations. They found the
source preprint and adjacent mixed-unitarity work, but no statement of
Theorem \ref{thm:main}. This is evidence, not a guarantee, of novelty.

Human review should focus on Lemma \ref{lem:semi}. Carath\'eodory gives a
uniform finite number of auxiliary unitary matrices, so Tarski--Seidenberg
applies in a fixed finite-dimensional real space. After that step, the
one-variable analytic argument is direct. The packet should remain a partial
result until the existence or impossibility of \(t_0<t_1\) is resolved.

\begin{thebibliography}{9}

\bibitem{BD}
B. V. Rajarama Bhat and R. Devendra,
\emph{On regions of mixed unitarity for semigroups of unital quantum
channels}, arXiv:2512.23598, version of 9 February 2026.
\href{https://arxiv.org/abs/2512.23598}{arXiv:2512.23598}.

\bibitem{BPR}
S. Basu, R. Pollack, and M.-F. Roy,
\emph{Algorithms in Real Algebraic Geometry}, second edition,
Algorithms and Computation in Mathematics 10, Springer, 2006.

\end{thebibliography}

\end{document}

